\chapter{hydrate rim measurements}

Belite shows a significant slower hydration, which is also visible in the hydrate rim growth (Figure \ref{fig:hydration-rim-results}).
After 28 days, Belite still shows smaller hydration rims as alite after 7 days.

\begin{figure}[h!]
    \centering

	\begin{subfigure}{0.45\textwidth}
		\includegraphics[width=\textwidth]{hydrate-rim/alite_result.pdf}
		\caption{Hydration of alite from 1 to 28 days. [Kleiner.2024]}
		\label{fig:hyd-rim-alite}
	\end{subfigure}
	\hfill
	\begin{subfigure}{0.45\textwidth}
		\includegraphics[width=\textwidth]{hydrate-rim/belite_result.pdf}
		\caption{Hydration development of belite from 1 to 28 days.}
		\label{fig:hyd-rim-belite}
	\end{subfigure}
    \caption{Comparison between the hydration progression measured using XRD (dashed lines) with $l_\text{max}$ (solid line).
	The pink markers (dotted line) indicate the amount of hydrate phase $\sigma_\text{hyd}$ estimated from the image segmentation. The red translucent areas indicate the confidence interval at a confidence level of where 33.3\,\% (dark) and 50\,\% (light). \label{fig:comp_hydration}}
    \label{fig:hydration-rim-results}
\end{figure}

\begin{itemize}
	\item C-S-H structure
	\item hadley-grain development more pronounced in belite
	\item gap between rim and clinker in belite
	\item TODO MIXTURES?
\end{itemize}

